\begin{table}[H]
\centering
\scriptsize
\caption{Composition of the control groups under observed and counterfactual connectivity}
\label{tab:direct_cf_composition}
\begin{tabular}{lcccccc}
\toprule
Control group & \shortstack{Control\\firms} & \shortstack{Log\\employment} & \shortstack{Employment} & \shortstack{Worker\\flows} & \shortstack{Flow\\partners} & \shortstack{Share in\\observed pure} \\
\midrule
All controls (Panel C) & 4,196 & 3.51 & 151.4 & 0.2413 & 6.00 & 0.460 \\
Observed pure (Panel A) & 1,931 & 2.38 & 22.4 & 0.2283 & 1.58 & 1.000 \\
Procedure A (cf. zero conn.) & 1,199 & 2.21 & 19.4 & 0.1766 & 0.43 & 0.926 \\
Procedure B (cf. bottom 46%) & 1,931 & 3.28 & 163.0 & 0.1797 & 3.29 & 0.623 \\
\bottomrule
\end{tabular}
\begin{minipage}{\textwidth}\vspace{2mm}\scriptsize
\textit{Notes:} This table compares the pre-treatment composition of the control groups entering the baseline direct-effect specification (all controls), the observed pure-control specification (untreated firms with zero observed connectivity to the treated set), and the two counterfactual-connectivity placebo selections, averaged across permutation draws. ``Flow partners'' is the number of analysis-sample firms with which the control firm exchanges any worker over 2007--2011. Procedure A keeps controls with counterfactual zero connectivity; Procedure B keeps the bottom 46.0\% of controls by counterfactual connectivity. The comparison shows that neither placebo rule reproduces the composition of the observed pure-control group, which should be borne in mind when reading the placebo distributions.
\end{minipage}
\end{table}
